% MỞ ĐẦU (LỜI MỞ ĐẦU - Chương 0)

\chapter*{MỞ ĐẦU} % Tên của chương
\addcontentsline{toc}{chapter}{MỞ ĐẦU} % Thêm tên chương vào mục lục

\label{Chapter0} % Để trích dẫn chương này ở chỗ nào đó trong bài, hãy sử dụng lệnh \ref{Chapter0} 

%----------------------------------------------------------------------------------------

Trong những năm gần đây, Hệ thống giao thông thông minh (Intelligent Transportation Systems - ITS) đã trở thành một phần không thể thiếu trong việc điều tiết và quản lý cơ sở hạ tầng. Trong đó, hệ thống camera giám sát đóng vai trò quan trọng giúp thu thập và phân tích dữ liệu lưu lượng phương tiện.

Tuy nhiên, trong môi trường thực tế, các hệ thống camera giao thông phải đối mặt với những rào cản vật lý vô cùng khắc nghiệt. Thực tế cho thấy, các yếu tố môi trường như sương mù, mưa lớn, hay điều kiện thiếu sáng vào ban đêm làm giảm độ tương phản, làm mờ và gây nhiễu ảnh nghiêm trọng. Điều này dẫn đến tình trạng các thuật toán nhận diện truyền thống dễ bị mất dấu đối tượng, nhận diện sai hoặc bỏ sót phương tiện, làm sai lệch kết quả thống kê lưu lượng giao thông.

Xuất phát từ nhu cầu thực tiễn đó, em quyết định thực hiện đề tài: \textbf{"Nhận diện và đếm số lượng phương tiện giao thông trong các điều kiện thời tiết khắc nghiệt"}.

Để đảm bảo tính khoa học, về mặt dữ liệu huấn luyện, đề tài sử dụng bộ dữ liệu TSBOW (Traffic Surveillance Benchmark for Occluded Vehicles Under Various Weather Conditions) chuyên biệt về giao thông trong điều kiện thời tiết xấu và bị che khuất. Về phương pháp triển khai mô hình, đề tài sử dụng mô hình mạng học sâu YOLO11 tích hợp thuật toán theo dõi đa đối tượng ByteTrack, xây dựng logic đếm xe để chạy thử nghiệm trên các video giao thông thực tế.

Để làm rõ hơn về đề tài này, bài báo cáo được chia làm các nội dung như sau:

\textbf{Chương 1: Tổng quan}

\textbf{Chương 2: Cơ sở lý thuyết}

\textbf{Chương 3: Thực nghiệm}

\textbf{Chương 4: Kết quả và đánh giá}

\textbf{Chương 5: Kết luận và hướng phát triển}
